
\section{Image directe}

Dans la suite il faut prendre la limite sur l'ultracatégorie slice $\bb_x(\A)$ et non pas sur la catégorie, \ie la limite donne des $\xi^h_s$ qui sont aussi stable par l'ultraprod de cet ultracat. (sinon la formule ne marche pas dans le cas topologique meme)

-----


On prend une map de petites vult bornées $f : \A\to\B$, sur les faisceau l'image inverse est donnée par la precomposition par $f$, on va donner une expression explicite de l'image directe (adjoint à droite de la précomp par $f$ donc v·u·right Kan extension le long de $f$).

On part d'un faisceau $X\in\sh(\A)$, et on construit $Y\in\sh(\B)$ de sorte à avoir $Y = f_*(X)$. Sur les objets on a 

\[Y(x) = \lim_{x\ultrato(f(a_s))_{s:\mu}}\int_{s:\mu}X(a_s).\]

Explicitement, la limite est prise le long de la slice catégorie "$x/_{\bb}f$" où une map de $x\ultrato(f(a_s))_{s:\mu}$ dans $x\ultrato(f(b_t))_{t:\nu}$ est donnée par des $(h_s : a_s\ultrato (c_{s,t})_{t:\nu_s})_{s:\mu}$ ainsi que d'un $\alpha : \nu\to\sum_{s:\mu}\nu_s$ tels que $x\ultrato(f(a_s))\stackrel{f(h_s)}\ultrato((f(b_{s,t})))$ s'épaissie le long de $\alpha$ en $x\ultrato(f(b_t))_{t:\nu}$.

Donc $Y(x)$ est la donnée de pour tout $h : x\ultrato (f(a_s))_{s:\mu}$, d'un $(\xi^h_s:X(a_s))_{s:\mu}$ qui sont stables par composition et épaississement. (peut-on voir $Y(x)$ plus simplement?)

Pour une ultraflèche $h : x\ultrato(y_s)$ dans $\B$, on obtient par precomposition par $h$

\[Y(x) = \lim_{x\ultrato(f(a_t))_{t:\nu}}\int_{t:\nu}X(a_t)\longrightarrow \lim_{(y_s\ultrato(f(a^s_{t}))_{t:\nu_s})_{s:\mu}}\int_{s:\mu}\int_{t:\nu_s}X(a^s_{t})\]

mais on a une injection 

\[\int_{s:\mu}Y(y_s) = \int_{s:\mu}\left(\lim_{y_s\ultrato(f(a^s_{t}))_{t:\nu_s}}\int_{t:\nu_s}X(a^s_{t})\right) \longhookrightarrow \lim_{(y_s\ultrato(f(a^s_{t}))_{t:\nu_s})_{s:\mu}}\int_{s:\mu}\int_{t:\nu_s}X(a^s_{t})\]

%mapping some $(((\xi(s,g_s,t):X(a^s_t))_{t:\nu_s})_{g_s:y_s\ultrato(f(a^s_{t}))_{t:\nu_s}})_{s:\mu}$ onto $(((\xi(s,g_s,t))_{t:\nu_s})_{s:\mu})_{(g_s :y_s\ultrato(f(a^s_{t}))_{t:\nu_s})_{s:\mu}}$

(en générale dans Set on a toujours une injection $\int\lim(-)\hookrightarrow\lim\int(-)$, cela se généralise aux autres ultracat?)

Il reste donc a montrer que la première equation se factorise par la seconde. 

PARTIE TECHNIQUE! (j'ai une preuve dégeu qui revient de partir d'un élément de l'ensemble de droite, de bien ordonner les $\{y_s\ultrato(f(a^s_t))\}$ pour tout $s$, puis de les regrouper en parallèle et d'utiliser l'élément du codom pour donner des valeurs aux $\int X(a^s_t)$, puis de montrer que le choix du parallèlisme était arbitraire) TROUVER MIEUX!

La fonctorialité de $Y$ se vérifie facilement, on obtient donc un faisceau sur $\B$, et il est clair que les homsets $\sh(\B)(Z, Y)$ et $\sh(\A)(f^*(Z), X)$ sont en correspondance naturelle pour $Z\in\sh(\B)$.



------



Il me semble que la formule ne marche pas... si on fixe deux ensembles $I$ et $A$.  On prend pour $\B$ une vult avec objets $p,(q_s)_{s:S},a$ et générée par $f_s^i : q_s\to a$ pour tout $i:I$ et $s:S$ et une ultraflèche $g : p\ultrato(q_s)_{s:\mu}$. On a donc $I^{\mu}$ ultraflèches $(f_s^{i_s})_{s:\mu}\circ g : p \ultrato a$. On prend $\A = \{a\}$ et le faisceau constant en $A$ dessus. Quand on étend avec la formule plus haut; la valeur en $p$ est donnée par $(A^{\mu})^{I^{\mu}} = [I^{\mu},A^{\mu}]$, et la valeur en un $q_s$ est donnée par $A^I$. Donc la fonctorialité donne une map $[I^{\mu},A^{\mu}] \to [I,A]^{\mu}$ et on peut vérifier (?) que ca devrait être l'inverse de l'injection canonique $[I,A]^{\mu} \hookrightarrow [I^{\mu},A^{\mu}]$, mais cette injection est stricte en générale (?).

Par exemple, pour $A = 2$, on a la valeur au dessus de $p$ est $P(I^{\mu})$, on prend $C\in P(I^{\mu})$ la partie des ultrafamilles constantes, on obtiendrait alors par fonctorialité des ultrafamilles $(J_s\subseteq I)_{s:\mu}$ (en regardant au dessus des $(q_s)$) tel que $(i_s)_{s:\mu}$ est constante ssi $\forall s:\mu, i_s\in J_s$, ce qui est impossible.

\section{Ultraidempotents}

To show that this is sufficient we will study more the notion of ultraretraction, and show that an ultraretract satisfies a (actually two) universal propertie(s).

Similarly to the usual notion of retract, we define the notion of ultraidempotent and give universal properties for the ultraretract of an ultraidempotent.

\begin{defn}
    Let $\X$ be any vultracat
    \begin{enumerate}
        \item an \emph{ultraidempotent} on an ultrafamily of objects $(a_s)_{s:\mu}$ is given by an ultrafamily of ultaarows $(\sigma_{s} : a_s \ultrato (a_t)_{t:\mu})_{s:\mu}$ satisfying for $s:\mu$ the following equation 
        \[(\sigma_t)_{t:\mu}\circ\sigma_s = \pi_2^*(\sigma_s)\]
        where
        \[(\sigma_t)_{t:\mu}\circ\sigma_s  :
            a_s\ultrato(a_t)_{t:\mu} \ultrato(a_u)_{(t:\mu)\otimes(u:\mu)}\]
        \[\sigma_s : a_{s}\ultrato(a_u)_{u:\mu}\]
        \[\pi_2 : (t:\mu)\otimes(u:\mu) \to (u:\mu)\]
        \item a spliting of such an ultraidempotent is given by an object $p$ and ultraarows $i : p\ultrato(a)_{s:\mu}$ and $(r_s : a_s\ultrato(p)_{\nu_s})_{s:\mu}$ such that 
        \[(i)_{\nu_s}\circ r_s = \pi_2^*(\sigma_s) \text{ and } (r_s)_{s:\mu}\circ i = \ !^*(\id_{p})\]
    \end{enumerate}
\end{defn}

\begin{prop}
    If $p$ is an ultraretract of $(a_s)_{s:\mu}$ with $i : p \ultrato (a_s)_{s:\mu}$ and $(r_s : a_s\ultrato(p)_{\nu})$ (on the same $\nu$) then we have an ultraidempotent on $((a_s)_{s:\mu})_{\nu}$
    \[(\sigma_{t,s} : a_s \ultrato (p)_{\nu} \ultrato ((a_u)_{u:\mu})_{\nu})_{(t,s):\nu\otimes\mu}\]
    moreover, $p$ is an spliting of this ultraidempotent.
\end{prop}


We compute the ultraidempotent in the examples above.
\begin{ex}
    
\end{ex}

\begin{prop}
    We assume that the thickening functors of $\X$ preserve the coequalizers of the form 
    \[(s:\mu)\otimes\lambda_s\otimes\lambda_s\otimes\nu_s\rightrightarrows(s:\mu)\otimes\lambda_s\otimes\nu_s\to(s:\mu)\otimes\nu_s\]
    (this is the case in particular if $\X$ locally sober)
    
    There is a natural bijection (in $q\in\X_0^{\op}$) between arrows $f : q\to p$ and ultraarows $g : q\ultrato(a_s)_{s:\mu}$ that are stable by postcomposition with the $\sigma_s$,
    \[(\sigma_s)_{s:\mu}\circ g = \pi_2^*(g)\]
    \ie such that
    \[q\stackrel{g}\ultrato(a_s)_{s:\mu}\stackrel{(\sigma_s)}\ultrato((a_t)_{t:\mu})_{s:\mu} = (a_t)_{(s,t):\mu\otimes\mu}\]
    is the thickening along $\pi_2 : \mu\otimes\mu\to\mu$ of
    \[q\stackrel{g}\ultrato(a_t)_{t:\mu}\]
\end{prop}
\begin{proof}
    From some $f : q\to p$ we get by postcomposition with the $i$ an ultraarow $g = i\circ f$
    \[g : p\to q \ultrato (a_s)_{s:\mu},\]
    and we need to show that this $g$ is stable by the $\sigma_s$
    \ie we want to show
    \[(\sigma_s)\circ i \circ f\ = \pi_2^*(i\circ f) : q\ultrato((a_t)_{t:\mu})_{s:\mu}\]
    and by \colorier{the hypothesis } it is enough to show that theses two map have the same thickening along 
    $\pi^{(s:\mu).\nu_s.\mu}_{\mu.\mu} : (s:\mu)\otimes\nu_s\otimes\mu \to \mu\otimes\mu$, that we can deduce from the following diagram:
    % https://q.uiver.app/#q=WzAsNixbMCwwLCJxIl0sWzIsMCwicCJdLFs0LDAsIihhX3MpX3tzOlxcbXV9Il0sWzUsMiwiKCgoYV90KV97dDpcXG11fSlfe1xcbnVfc30pX3tzOlxcbXV9Il0sWzMsMiwiKChwKV97XFxudV9zfSlfe3M6XFxtdX0iXSxbMSwyLCIoKHEpX3tcXG51X3N9KV97czpcXG11fSJdLFswLDEsImYiXSxbMSwyLCJpIl0sWzAsMiwiZyIsMCx7ImN1cnZlIjotNSwic3R5bGUiOnsiYm9keSI6eyJuYW1lIjoiZGFzaGVkIn19fV0sWzIsMywiKFxccGlfMl4qKFxcc2lnbWFfcykpX3tzOlxcbXV9IiwxLHsic3R5bGUiOnsiYm9keSI6eyJuYW1lIjoiZGFzaGVkIn19fV0sWzIsNCwiKHJfcylfe3M6XFxtdX0iLDEseyJzdHlsZSI6eyJib2R5Ijp7Im5hbWUiOiJkYXNoZWQifX19XSxbNCwzLCIoKGkpX3tcXG51X3N9KV97czpcXG11fSIsMl0sWzEsNCwiXFxkZWx0YSIsMix7InN0eWxlIjp7ImJvZHkiOnsibmFtZSI6ImRhc2hlZCJ9fX1dLFswLDUsIlxcZGVsdGEiLDIseyJzdHlsZSI6eyJib2R5Ijp7Im5hbWUiOiJkYXNoZWQifX19XSxbNSw0LCIoKGYpX3tcXG51X3N9KV97czpcXG11fSIsMl0sWzUsMywiKChnX3MpX3tcXG51X3N9KV97czpcXG11fSIsMix7ImN1cnZlIjo1LCJzdHlsZSI6eyJib2R5Ijp7Im5hbWUiOiJkYXNoZWQifX19XV0=
\[\begin{tikzcd}
	q && p && {(a_s)_{s:\mu}} \\
	\\
	& {((q)_{\nu_s})_{s:\mu}} && {((p)_{\nu_s})_{s:\mu}} && {(((a_t)_{t:\mu})_{\nu_s})_{s:\mu}}
	\arrow["f", from=1-1, to=1-3]
	\arrow["g", curve={height=-30pt}, dashed, from=1-1, to=1-5]
	\arrow["\delta"', dashed, from=1-1, to=3-2]
	\arrow["i", from=1-3, to=1-5]
	\arrow["\delta"', dashed, from=1-3, to=3-4]
	\arrow["{(r_s)_{s:\mu}}"{description}, dashed, from=1-5, to=3-4]
	\arrow["{(\pi_2^*(\sigma_s))_{s:\mu}}"{description}, dashed, from=1-5, to=3-6]
	\arrow["{((f)_{\nu_s})_{s:\mu}}"', from=3-2, to=3-4]
	\arrow["{((g)_{\nu_s})_{s:\mu}}"', curve={height=30pt}, dashed, from=3-2, to=3-6]
	\arrow["{((i)_{\nu_s})_{s:\mu}}"', from=3-4, to=3-6]
\end{tikzcd}\]
    and 
    \[(\pi_2^*(\sigma_s)) \circ g = ((s:\mu)\otimes\pi_2)^*((\sigma_s)\circ g) = (\pi^{(s:\mu).\nu_s.\mu}_{\mu.\mu})^*((\sigma_s)\circ g)\]
    \[(g)_{\nu}\circ \delta= (g)_{\nu}\circ (\pi^{\nu}_{\ptuf})^*\id_q = (\pi^{\nu}_{\ptuf}\otimes\mu)^*(g) = (\pi^{(s:\mu).\nu_s.\mu}_{\mu.\mu})^*(\pi_2^*(g))\]

    \

    Conversely, from some $g:q\ultrato(a_s)_{s:\mu}$ stable under the $(\sigma_s)$, we can postcompose with the $(r_s)$ to get an ultraarow $f' = (r_s)_{s:\mu}\circ g$
    \[f' : q\ultrato (a_s)_{s:\mu}\ultrato((p)_{\nu_s})_{s:\mu}\]
    and we need to show that this $f'$ can be unthicken into some $f : q\to p$.
% https://q.uiver.app/#q=WzAsNixbMCwwLCJxIl0sWzIsMCwiKGFfcylfe3M6XFxtdX0iXSxbNCwwLCIoKHApX3tcXG51X3N9KV97czpcXG11fSJdLFsyLDIsIigoKGFfdClfe3Q6XFxtdX0pX3tcXG51X3N9KV97czpcXG11fSJdLFs0LDIsIigoKChwKV97XFxudV90fSlfe3Q6XFxtdX0pX3tcXG51X3N9KV97czpcXG11fSJdLFswLDVdLFswLDEsImciXSxbMSwyLCIocl9zKV97czpcXG11fSJdLFsyLDMsIihpKSIsMV0sWzMsNCwiKHJfdCkiLDJdLFsyLDQsIigoXFxkZWx0YSlfe1xcbnVfc30pX3tzOlxcbXV9Il0sWzEsMywiKFxccGlfMl4qKFxcc2lnbWFfcykpX3tzOlxcbXV9IiwxXSxbMCwzLCJcXGFscGhhXiooKFxcc2lnbWFfcylcXGNpcmMgZykgIiwxXSxbMCwzLCJcXGFscGhhXiooXFxwaV8yXiooZykpIiwyLHsiY3VydmUiOjV9XSxbMCwyLCJmIiwwLHsiY3VydmUiOi01fV0sWzEyLDEzLCIiLDEseyJzaG9ydGVuIjp7InNvdXJjZSI6NDAsInRhcmdldCI6NDB9LCJzdHlsZSI6eyJoZWFkIjp7Im5hbWUiOiJub25lIn19fV1d
\[\begin{tikzcd}
	q && {(a_s)_{s:\mu}} && {((p)_{\nu_s})_{s:\mu}} \\
	\\
	&& {(((a_t)_{t:\mu})_{\nu_s})_{s:\mu}} && {((((p)_{\nu_t})_{t:\mu})_{\nu_s})_{s:\mu}} 
	\arrow["g", from=1-1, to=1-3]
	\arrow["f'", curve={height=-30pt}, from=1-1, to=1-5]
	\arrow[""{name=1, anchor=center, inner sep=0}, "{(\pi^{(s:\mu).\nu_s.\mu}_{\mu.\mu})^*((\sigma_s)\circ g)}"', from=1-1, to=3-3]
	\arrow["{(r_s)_{s:\mu}}", from=1-3, to=1-5]
	\arrow["{(\pi_2^*(\sigma_s))_{s:\mu}}"{description}, from=1-3, to=3-3]
	\arrow["{(i)}"{description}, from=1-5, to=3-3]
	\arrow["{((\delta)_{\nu_s})_{s:\mu}}", from=1-5, to=3-5]
	\arrow["{(r_t)}"', from=3-3, to=3-5]
\end{tikzcd}\]
The upper way gives
\[(\delta)_{\nu}\circ f' = ((\pi^{\nu}_{\ptuf})^*(\id_p))_{\nu} \circ f' =  (\nu\otimes\pi^{\nu}_{\ptuf})^*(f') = \pi_1^*(f')\text{ , for } \pi_1 : \nu\otimes\nu\to\nu\]
and the downward way gives
\[((r_t)_{t:\mu})_{\nu}\circ(\pi^{(s:\mu).\nu_s.\mu}_{\mu.\mu})^*((\sigma_s)\circ g) = ((r_t)_{t:\mu})_{\nu}\circ(\pi^{(s:\mu).\nu_s.\mu}_{\mu.\mu})^*(\pi_2^*(g))\]
\[= ((r_t)_{t:\mu})_{\nu}\circ(\pi^{\nu.\mu}_{\mu})^*(g) = (\pi^{\nu.(t:\mu).\nu_t}_{(t:\mu).\nu_t})^*((r_t)_{t:\mu}\circ g) = \pi_2^*(f')\text{ , for } \pi_2 : \nu\otimes\nu\to\nu\]
and so \colorier{by the hypothesis } $f'$ can be unthicken in some (unique) $f : q\to p$. 

    \

    We now show that these establishes a bijective correspondance. From some $f:p\to q$ we first get $i\circ f : q \ultrato(a_s)_{s:\mu}$ stable under the $(\sigma_s)$ and then $f' = (r_s)_{s:\mu}\circ i\circ f : q\ultrato((p)_{\nu_s})_{s:\mu}$, and as $(r_s)_{s:\mu}\circ i$ is the diagonal, $f'$ is indeed the thickening of $f$.

    Conversely, from some $g:q\ultrato(a_s)_{s:\mu}$ stable under the $(\sigma_s)$ we have that $(r_s)_{s:\mu}\circ g : q\ultrato ((p)_{\nu_s})_{s:\mu}$ is the thickening of some $f:q\to p$, and $g' = i\circ f : p \ultrato (a_s)_{s:\mu}$ satisfied so
    \[(\pi^{\nu.\mu}_{\mu})^*(g') = (\pi^{\nu.\mu}_{\mu})^*(i\circ f) = (i)_{\nu} \circ (!_{\nu})^*(f) = (i)_{\nu} \circ (r_s)_{s:\mu}\circ g = ((i)_{\nu_s}\circ r_s)_{s:\mu} \circ g \]
    \[= ((\pi^{\nu_s.\mu}_{\mu})^*(\sigma_s))_{s:\mu}\circ g = (\pi^{\nu.\mu}_{\mu.\mu})^*((\sigma_s)_{s:\mu}\circ g) = (\pi^{\nu.\mu}_{\mu.\mu})^*((\pi^{\mu.\mu}_{\ptuf.\mu})^*(g)) = (\pi^{\nu.\mu}_{\mu})^*(g)\]
    and so \colorier{using the hypothesis}, we get $g=g'$.
    
\end{proof}

We can analogously give a description of ultraarows $q\ultrato (p)_{\lambda}$ and ultraarows $p\ultrato(q_l)_{l:\lambda}$




\section{The tautness condition}

The notion of v-ult is suppose to categorify the one of topological spaces. But as noticed in \cite{G,JUS,Ali} a v-ultraposet, even small, is not always the point of a topological space, it misses one condition! 

\begin{prop}
    Let $\X$ be a small v-ultraposet, then the following are equivalent
    \begin{enumerate}
        \item $\X = \pt(T)$ for some topological space $T$.
        \item $\X$ is locally sober.
        \item for any $\alpha : \nu\to\mu$ in $\UF$, if there is an ultraarow $a\ultrato(b_{f(t)})_{t:\nu}$ in $\X$ then there is an ultraarow $a\ultrato(b_s)_{s:\mu}$.
    \end{enumerate}
\end{prop}

This last condition prescribes the thickening operation on $\X$, in general the thickening operation can give a lot of complexity to the notion of v-ult. The idea of the following is to prescribes its behavior by categorifying the condition \ref{}.

\begin{defn}
    For $a$ and $(b_s)_{s:\mu}$ points of a vult $\X$, the thickening functor is defined to be
    \begin{align*}
        (\UF/{\mu})^{\op} & \to \Set \\
        (\nu\stackrel{\alpha}\to\mu) & \mapsto \X(a,(b_{\alpha(t)})_{t:\nu})
    \end{align*}
\end{defn}

We want to understand the behavior of the thickening functor in (locally) sober v-ultracat.

\begin{lem}
    For any set $A$ and ultrafilter $\mu$, the following is an equalizer of sets
    \[A \to A^{\mu}\rightrightarrows A^{\mu\otimes\mu}\]
    We deduce that the following is a coequalizer in $\UF$
    \[\mu\otimes\mu\rightrightarrows\mu\to\ptuf\]
\end{lem}
\begin{proof}
    Let $(a_s)_{s:\mu}$ be such that $(\forall s:\mu)(\forall t:\mu)\ a_s = a_t$, then we can find a $t_0$ such that $(\forall s:\mu)\ a_s = a_{t_0}$, \ie $(a_s)_{s:\mu}$ is almost constant as wanted.

    Now, let $\alpha : \mu\to\nu$ be such that $\alpha\circ\pi_1 = \alpha\circ\pi_2$, and we want to show that $\nu = \ptuf$. Let $T$ be a defining set of $\nu$ so that $(\alpha_l)\in T^{\mu}$, the hypothesis says that $(\alpha_l)$ belongs to the equalizer of $T^{\mu}\rightrightarrows T^{\mu\otimes\mu}$, and so to $T$ \ie it is constant.
\end{proof}

\begin{prop}
    The diagram 
    \[(s:\mu)\otimes\lambda_s\otimes\lambda_s\otimes\nu_s\rightrightarrows(s:\mu)\otimes\lambda_s\otimes\nu_s\to(s:\mu)\otimes\ptuf\otimes\nu_s\]
    is a coequalizer in $\UF/\nu$ (with $\nu \coloneqq \sum_{s:\mu}\nu_s$).
    
    Moreover, in a sober vult, the thickening functor $(\UF/\nu)^{\op} \to \Set$ sends it to a equalizer. 
\end{prop}
\begin{proof}
    First in the case $\lambda\otimes\lambda\rightrightarrows\lambda\to\ptuf$. Then we know it is a coeq in $\UF$ and we need to show that $\mu\mapsto\Set(A,B^{\mu})$ sends it to a eq, as $\Set(A,-)$ is continuous it is enough to show that it is the case for $\mu\mapsto B^{\mu}$ and this is exactly the above lemma.

    ...
    
\end{proof}

\begin{defn}
    A v-ultracategory $\X$ is said \emph{taut} if all the thickening functors preserves these coequalizers,
    \ie for any $a$ and $((b_{s,t})_{t:\nu})_{s:\mu}$ the following is an equalizer of sets
    \[\X(a,(b_{s,t})_{(s:\mu)\otimes(t:\nu_s)})\to\X(a,(b_{s,t})_{(s:\mu)\otimes\lambda_s\otimes(t:\nu_s)})\rightrightarrows\X(a,(b_{s,t})_{(s:\mu)\otimes\lambda_s\otimes\lambda_s\otimes(t:\nu_s)})\]
    \ie an ultraarow $f : a\ultrato (b_{s,t})_{(s:\mu)\otimes\lambda_s\otimes(t:\nu_s)}$ unthicken to an ultraarow $a\ultrato(b_{s,t})_{(s:\mu)\otimes(t:\nu_s)}$ iff the two thickening to $a\ultrato (b_{s,t})_{(s:\mu)\otimes\lambda_s\otimes\lambda_s\otimes(t:\nu_s)}$ are identic, and then such a dethickening is unique.
\end{defn}

\begin{prop}
    Locally sober v-ults are taut.
\end{prop}

\begin{prop}
    If $\X$ is taut, then the thickening functors $(\UF/\mu)^{\op}\to\Set$ are sheaves for the atomic topology.
\end{prop}
\begin{proof}
    Let $\alpha : \nu\to\mu$ and $f : a\ultrato(b_{\alpha(t)})_{t:\nu}$. We need to show that if for any $\beta,\gamma:\lambda\rightrightarrows\nu$ equalizing $\alpha$ the thickening of $f$ along $\beta$ and $\gamma$ are equals, then $f$ unthicken to some unique $a\ultrato(b_s)_{s:\mu}$.

    By the ‘choice lemma’ in $\UF$ \cite{SUJ} we can ‘make $\alpha$ a projection’ \ie we can find $\kappa$ and $\beta : \kappa\otimes\mu \to \nu$ such that $\alpha\circ\beta = \pi_1$. Then both maps $\kappa\otimes\kappa\otimes\mu\rightrightarrows\kappa\otimes\mu\to\nu$ equalize $\alpha$, and so $\beta^*f : a \ultrato(b_s)_{\kappa\otimes\mu}$ have the same thickening along both projections $\kappa\otimes\kappa\otimes\mu$, and so by hypothesis it unticken into a unique $a\ultrato(b_s)_{s:\mu}$.
\end{proof}

\begin{rmk}
    Is this an iff?? Can we restrict the taut condition to some smaller class of coeq (only the one used in the proof above, i think yes)?? Does tautness implies the preservation of all conn colimits (at least of pushout built with $\otimes$ seems yes)?? 
\end{rmk}

\begin{prop}
    It is actually an iff!!
\end{prop}
\begin{proof}
    Let's do the easy case, suppose we have $A : \UF^{\op}\to\Set$ sheaf, and we want to show that $\mu\otimes\mu\rightrightarrows\mu\to\ptuf$ is sent to an equalizer by $A$. Let $x : A(\mu)$ such that $\pi_1^*(x) = \pi_2^*(x)$, to use the sheaf condition it is enough to show that for any $f,g:\nu\rightrightarrows\mu$ we have $f^*(a) = g^*(a)$.

    A REVOIR! \colorier{We use the choice lemma twice to get $h : \lambda\otimes\mu\otimes\mu\to\mu$ such that $f\circ h = \pi^{\lambda.\mu.\mu}_{\ptuf.\mu.\ptuf}$ and $g\circ h = \pi^{\lambda.\mu.\mu}_{\ptuf.\ptuf.\mu}$.}

    Then $h^*(f^*(x)) = (\pi^{\lambda.\mu.\mu}_{\ptuf.\mu.\ptuf})^*(x) = (\pi^{\lambda.\mu.\mu}_{\mu.\mu})^*(\pi_1^*(x)) = (\pi^{\lambda.\mu.\mu}_{\mu.\mu})^*(\pi_2^*(x)) = (\pi^{\lambda.\mu.\mu}_{\ptuf.\ptuf.\mu})^*(x) = h^*(g^*(x))$.

    So $h^*(f^*(x)) = h^*(g^*(x))$ and only remains to show that $h^*$ is injective, and this follows from the sheaf condition.
\end{proof}

\begin{proof}
    On montre qu'un atomic-sheaf respect le coeq $\sigma\otimes\sigma\rightrightarrows\sigma\to\pt$. C'est le seul qu'on a besoin dans la section sur les ultra-idem (vrai?), mais j'espère qu'une preuve similaire s'adapte pour les autres coeq basé sur le otimes.

    On prend $F\in\Sh(\UF_{at})$ et $\xi \in F(\sigma)$ tq $\pi_1^*(\xi) = \pi_2^*(\xi)\in F(\sigma\otimes\sigma)$. On prend $\phi,\psi:\tau\rightrightarrows\sigma$ and we need to show that $\phi^*(\xi) = \psi^*(\xi)$.
\[\begin{tikzcd}
	{\tau\otimes\tau} & \tau \\
	{\sigma\otimes\sigma} & \sigma
	\arrow["{\pi_1}"{description}, shift left=1, from=1-1, to=1-2]
	\arrow["{\pi_2}"{description}, shift right=1, from=1-1, to=1-2]
	\arrow["{\phi\otimes\psi}"', from=1-1, to=2-1]
	\arrow["\phi"', shift right, from=1-2, to=2-2]
	\arrow["\psi", shift left, from=1-2, to=2-2]
	\arrow["{\pi_2}"{description}, shift right=1, from=2-1, to=2-2]
	\arrow["{\pi_1}"{description}, shift left=1, from=2-1, to=2-2]
\end{tikzcd}\]

As $F$ sends any map to a mono, it is enough to show that $\pi_1^*(\phi^*(\xi)) = \pi_1^*(\psi^*(\xi))$

\[\pi_1^*(\phi^*(\xi)) = (\phi\otimes\psi)^*(\pi_1^*\xi) = (\phi\otimes\psi)^*(\pi_2^*\xi) = \pi_2^*(\psi^*(\xi))\]

Hence, for all $\phi,\psi : \tau\rightrightarrows\sigma$ we have $\pi_1^*(\phi^*(\xi)) = \pi_2^*(\psi^*(\xi))$

In particular (using $\psi\otimes\psi$ instead of $\phi\otimes\psi$) we get 
$\pi_1^*(\psi^*(\xi)) = \pi_2^*(\psi^*(\xi))$.
and so 
$\pi_1^*(\phi^*(\xi)) = \pi_2^*(\psi^*(\xi)) = \pi_1^*(\psi^*(\xi))$
as wanted.

\end{proof}


------

We first notice that the point functor from $\TopSp$ to vultraposets have a left-adjoint.

\begin{defn}
    Let $A$ be a small v-ultraposet, we denote by $\topo{A}$ the v-ultraposet having the same object that $A$ and with an ultraarow ultraarow $a\ultrato(b_s)_{s:\mu}$ in $\topo(A)$ iff there is an ultraarow $a\ultrato(b_{f(t)})_{t:\nu}$ in $A$ for some $\alpha : \nu\to\mu$.
\end{defn}
\begin{prop}
    The vultraposet $\topo(A)$ is locally sober and it is universally so, \ie for a topological space $T$ the functor $\vUlt(\topo(A),T)\to\vUlt(A,T)$ is an equivalence. Hence $\topo$ gives a functor from vultraposet to $\TopSp$ that is left adjoint to the point functor.
\end{prop}

Motivated by the fact that (we expect?) small vultraposets to be dense into vultracategories, we can define a categorify version of the condition \ref{} by,

\begin{defn}
    A v-ultracategory $\X$ is said \emph{taut} if for all small vultraposet $A$, the precomposition functor $\vUlt(\topo(A),\X)\to\vUlt(A,\X)$ is an equivalence.
\end{defn}

hence, being ‘taut’ is a local soberness property, it means that for the point of view of a 0-dimensional vult the vult $\X$ is locally sober.

It happens that all example arising naturally are taut, even the non locally sober ones! for example the free ultracategories or the models of a continuous theory.

Whereas being locally sober is hard to describe algebraically, it is not the case of the tautness!

\begin{prop}
    A vult $\X$ is taut if anf only if for $a$ and $(b_s)_{s:\mu}$ points of $\X$, the thickening functor 
    \begin{align*}
        (\UF/{\mu})^{\op} & \to \Set \\
        (\nu\stackrel{\alpha}\to\mu) & \mapsto \X(a,(b_{\alpha(t)})_{t:\nu})
    \end{align*}
    is a sheaf for the atomic topology on $\UF/{\mu}$.
\end{prop}
\begin{proof}
    
\end{proof}

\begin{rmk}
    In general for $\X$ a sober vult, the thickening functor preserves all connected colimits. Asking for the thickening functor to preserves all connected colimits is à priori a stronger condition than the sheaf condition, we leave open the question to know if it is actually a stronger condition. It might be equivalent to the sheaf condition to ask the preservations of all connected colimits (or only to coeq/pushout?).
\end{rmk}


\section{Ultraretract, soberness, surjectivity}

Naive idea of retract, not the case for topological spaces, extracting ultraretract from Johnston proposition.

\begin{defn}
    Let $\X$ be any vultracat
    \begin{enumerate}
        \item an \emph{ultraidempotent} on an ultrafamily of objects $(a_s)_{s:\mu}$ is given by an ultrafamily of ultaarows $(\sigma_{s} : a_s \ultrato (a_t)_{t:\mu})_{s:\mu}$ satisfying for $s:\mu$ the following equation 
        \[(\sigma_t)_{t:\mu}\circ\sigma_s = \pi_2^*(\sigma_s)\]
        where
        \[(\sigma_t)_{t:\mu}\circ\sigma_s  :
            a_s\ultrato(a_t)_{t:\mu} \ultrato(a_u)_{(t:\mu)\otimes(u:\mu)}\]
        \[\sigma_s : a_{s}\ultrato(a_u)_{u:\mu}\]
        \[\pi_2 : (t:\mu)\otimes(u:\mu) \to (u:\mu)\]
        \item a spliting of such an ultraidempotent is given by an object $p$ and ultraarows $i : p\ultrato(a)_{s:\mu}$ and $(r_s : a_s\ultrato(p)_{\nu_s})_{s:\mu}$ such that 
        \[(i)_{\nu_s}\circ r_s = \pi_2^*(\sigma_s) \text{ and } (r_s)_{s:\mu}\circ i = \ !^*(\id_{p})\]
    \end{enumerate}
\end{defn}


\begin{lem}
    In $\Set$, an ultraretract is given by $i : P\hookrightarrow\int_{s:\mu}A_s$ with maps $r_s : A_s\to P^{\nu_s}$ such that $(r_s)\circ i = \delta$. Then $i$ is an injection whose image is given by the $(a_s)_{s:\mu}$ satisfying for $\mu$-all $s$
    \[i^{\nu_s}(r_s(a_s)) = ((a_u)_{u:\mu})_{\nu_s}.\]
\end{lem}
\begin{proof}
    As $(r_s)\circ i = \delta$ is an injection, so too is $i$. We now prove the characterization of the image. To simplify the notation, we denote $\nu \coloneqq \sum_{s:\mu}\nu_s$.
    
    Let $i(p) = (a_s)_{s:\mu}$ be in the image, then $(r_s)\circ i = \delta$ gives that $(r_s(a_s))_{s:\mu} = (p)_{\nu}$ \ie we have $r_s(a_s) = (p)_{\nu_s}$ for $s:\mu$, and so $i^{\nu_s}(r_s(a_s)) = (i(p))_{\nu_s} = ((a_u)_{u:\mu})_{\nu_s}$ as wanted.
    
    Conversely if $(a_s)_{s:\mu}$ satisfies the above condition, we first show that $(r_s(a_s))_{s:\mu}\in P^{\nu}$ is actually a constant ultrafamily: $(r_s(a_s))$ is maped by $i^{\nu}$ on the element $((a_u)_{u:\mu})_\nu \in (\int_{u:\mu}A_u)^{\nu}$ that belong to $\delta(\int_{u:\mu}A_u) \subseteq (\int_{u:\mu}A_u)^{\nu}$, and as $i$ is injective this implies that $(r_s(a_s))_{s:\mu}\in P^{\nu}$ should be in $\delta(P) \subseteq P^{\nu}$. Hence $(r_s(a_s))_{s:\mu}$ is of the form $(p)_{\nu}$ for some $p\in P$, and $(i(p))_{\nu_s} = i^{\nu_s}((p)_{\nu_s}) = i^{\nu_s}(r_s(a_s)) = ((a_u)_{u:\mu})_{\nu_s}$, and so $i(p) = (a_u)_{u:\mu}$ as wanted.
\end{proof}



Definition of ultraidem and ultraretracts, eg (retract, filtered colim, subclosure).

Computation of ultraretracts in Set, eso up to uret implies surj.

Discussion about the ultra Cauchy completion.




Correction of the fullness into full up to thick: two reason (factorization and to the counterex with hyperconnected)

Full utt implies hyperconnected. 

topological reflection = posetal reflection



\section{Characterizations of surjective and hyperconnected map}

Type of an ultrasheaf, and thm characterization of the image of an ultrasheaf.

Corolary giving other direction for surj and hc, and local soberness if faith in Set.


\section{Exactness properties of UF}

We denote $\UF$ the category of ultrafilters, $\eF$ the category of filters, and $\eG$ the category of germs of topological spaces (as considered by Moerdijk ‘‘A model for intuitionistic non-standard arithmetic’’).

\begin{defn}
    There are functors \[\UF\hookrightarrow\Fam(\UF)\hookrightarrow\eF\hookrightarrow\eG\]
    with right adjoints \[\Fam(\UF)\leftarrow\eF\leftarrow\eG\]
    and a left adjoint \[\eF\leftarrow\eG\]
    \begin{enumerate}
        \item the inclusion $\UF\hookrightarrow\Fam(\UF)$ is given by Yoneda.
        \item the inclusion $\Fam(\UF)\hookrightarrow\eF$ is defined by extending the above by noticing that $\eF$ has small coproducts : the coproduct of filters $(S_i,\F_i)_i$ is given by $(\sum S_i, \{\sum A_i \ |\ (A_i\in\F_i)_i\})$.
        \item the right adjoint $\F\to\Fam(\UF)$ sends a filter to the family of ultrafilers extending it.
        \item the functor $\F\hookrightarrow\G$ sends a filter $(S,\F)$ on the space $S\cup\infty$ with the topology given by $U\in S\cup\infty$ is open iff $(\infty\in U\implies U\cap S \in \F)$.
        \item the right adjoint $\G\to\F$ sends a germ $(T,x)$ to the filter of neighborhood of $x$ on $T_{\setminus\{x\}}$ (\ie $A \in T_{\setminus\{x\}}$ is large iff $A\cup\{x\}\subseteq T$ is a neighborhood of $x$).
        \item the right adjoint composite $\G\to\Fam(\UF)$ sends a germ $(T,x)$ to the set of ultrafilters $\mu\in\beta(T)$ converging to $x$.
        \item the left adjoint $\G\to\F$ sends a germ $(T,x)$ to the filter of neighborhood of $x$ in $T_{\setminus\{x\}}/\sim$ (where we quotient by the connected component of $T$, we use that locally-connected spaces is a reflective subcat).
    \end{enumerate}
\end{defn}

\begin{prop}
    NOOO, only finitely complete and coproducts, no coeq (see Blass)! -- $\G$ is cocomplete and complete (infinite products are given by partial dependent product!).
\end{prop}
\begin{prop}
    The product of $(T_i,x_i)$ is given by the space $\prod^p T_i$ of partial families, with the topology given by the basis $\{(t_i)\in\prod^p T_i \ | \ (t_{i_k}\in U_{i_k})\}$ for a finite number of $U_{i_k}\in T_{i_k}$. 
\end{prop}

\begin{cor}
    NOOO, actually $\UF$ has coeq, but not $\F$!! -- $\F$ and $\Fam(\UF)$ are bicomplete as coreflective sub-categories of a bicomplete one. So $\UF$ possesses multi-limits and connected coproducts. 
\end{cor}

We can give now an explicit description of (co)limits in $\F$ and $\UF$.

\begin{itemize}
    \item the initial object in $\F$ is the unique filter on the empty set (corresponding to the point in $\G$) and the terminal object is the unique filter on a singleton (corresponding to Sierpinski in $\G$).
    \item coproduct in $\F$ are given above, and products are computed as in $\G$, \ie we take the filter on the set of partial families, for $(S_i,\F_i)$ filters the product is $(\prod^p S_i,\prod\F_i)$
    \item multiproducts in $\UF$ are given by computing the product in $\F$ and taking all possible ultrafilter extensions of the result. connected colimits are computed as in $\F$, let look at the filtered colimit case.
    \item filitered colimit of a countable chain of filters $(S_0,\F_0)\to(S_1,\F_1)\to\cdots$ is given by...
\end{itemize}

\begin{defn}
    Let $(\mu_i, S_i)_{i:I}$ be a family of ultrafilters. We denote $S \coloneqq (i:I)\rightharpoonup S_i$ the set of partial dependent functions, and $\ev_i : S\rightharpoonup S_i$ the partial evaluation. For $A_{1}\in \mu_{i_1},\cdots,A_{n}\in \mu_{i_n}$ we denote $S_{A_1,\cdots,A_n}\coloneqq\{s:S \ | \ s(i_1)\in A_1\land \cdots\land s(i_n)\in A_n\}$. It is clear that these subsets are non-empty and closed by finite intersections, so they form a filter. Moreover, an ultrafilter $\mu\in\beta(S)$ extends this filter if and only if all $\dom(\ev_i) = \{s:S \ | \ i\in\dom(s)\}$ are $\mu$-large, and $(\ev_i)_*(\mu) = \mu_i$. 
\end{defn}
\begin{lem}
    There is a natural bijection $\sum_{\mu \text{ st } (\ev_i)_*(\mu) = \mu_i}\UF(-,\mu) \stackrel{\cong}\to \prod_{i:I}\UF(-,\mu_i)$ (this gives a description of the multi-products in $\UF$).
\end{lem}
\begin{proof}
    The direct map is given by postcomposing with the $\ev_i$. Suppose now we have $(\nu\to\mu_i)_{i:I}$, and so $(f_i : T\rightharpoonup S_i)_{i:I}$ all defined on $\nu$-large subsets of $T$. We define $f:T\to S$ by mapping $t:T$ on the partially defined $(f_i(t))_{i:I}\in S$, and it is clear that $f_i = \ev_i\circ f$ and so by letting $\mu\coloneqq f_*(\nu)\in\beta(S)$, we have a map $f\in\UF(\nu,\mu)$ with $(\ev_i)_*(\mu) = \mu_i$. And it is easy to see that this is well-defined up to iso of ultrafilters, and that that it is an inverse.
\end{proof}

\begin{prop}
    Let $\X\coloneqq\vpt(\E)$ the v-ult of points of a topos $\E$ and $X\subseteq\Ob(\X)$ a class of points of $E$. For $p\in\pt(\E)$ the following are equivalents:
    \begin{enumerate}
        \item $p$ belongs to $\E_{\res{X}}$, \ie $p$ is belongs to any subtopos of $\E$ containing $X$, \ie the restriction functor $\sh(\X_{\res{X\cup \{p\}}}) \to \sh(\X_{\res{X}})$ is an equivalence.
        \item $p$ is an ultraretract of an ultrafamily $(x_s)_{s:\mu}$ in $X$, \ie there exists $v : p \ultrato (x_s)_{s:\mu}$ and $(u_s : x_s\ultrato (p)_{\nu_s})_{s:\mu}$ ultraarrows, such that the composite $(u_s)\circ v : p \ultrato(p)_{\sum_{s:\mu}\nu_s}$ is the ultradiagonal.
    \end{enumerate}
\end{prop}

\begin{lem}%\label{Lemprebound}
    Let $\E$ be a topos, and $U : \E\to\SetO$ a localic morphism (\ie a prebound). The function $\vpt(\E)(p,(q_s)_{s:\mu})\to \Set(U_p,\int_{s:\mu}U_{q_s})$ is injective, and its image is given by the function $f : U_p\to\int_{s:\mu}U_q$ such that for any $V\subseteq U^n$, $f^n : U^n\to\int_{s:\mu}U^n_q$ restricts to $V\to\int_{s:\mu}V_q$.
\end{lem}
\begin{proof}
    Suppose there is such an $f$, we want to define $\alpha : \Nat_{E:\E}(E_p,\int_{s:\mu}E_{q_s})$. Let $E\in\E$, by the prebound hypothesis, $E$ is a quotient of some $V\subseteq U^n$. 
    
    Hence $E = \coeq(W\rightrightarrows V)$ with $W\coloneqq V\times_EV\subseteq V\times V\subseteq U^{2n}$. So
    \[E_p = \coeq(W_p\rightrightarrows V_p)\]
    and (as $\int_{s:\mu} : \Set^S\to\Set$ is a pretopos morphism).
    \[\int_{s:\mu}E_{q_s} = \coeq\left(\int_{s:\mu}W_{q_s}\rightrightarrows \int_{s:\mu}V_{q_s}\right)\]

    Using the hypothesis we get maps $f^{2n} : W_p \to \int W_{q_s}$ and $f^n : V_p \to \int V_{q_s}$ that induces a map $f_E : E_p \to \int_{s:\mu}E_{q_s}$ as wanted. We now need to check the naturality, let $\alpha : E\to E'\in\E$, we want to show that $(\int \alpha_{q_s}) \circ f_E = f_{E'}\circ\alpha_p$, by factorization it is enough to show it for $\alpha$ surjection and $\alpha$ mono. It is enough to show that we can chose a presentation of $E$ (resp. $E'$) by $V$ (resp. $V'$) with some map $V\hookrightarrow V'$ above $\alpha$: for $\alpha$ surjection, we chose a $V\twoheadrightarrow E$ presenting $E$ and this $V$ also present $E'$; for $\alpha$ mono, we chose a $V\twoheadrightarrow E'$ and pullback it to get a corresponding presentation of $E$.
\end{proof}


\begin{proof}
    We consider the indexing \colorier{(large)} set $S$ where an element $s:S$ is given by an ultraarow $a_s \stackrel{r_s}\ultrato (p)_{\nu_s}$ with $a_s\in A$ and for each $\xi:U_p$ a choice of a lift of $r_s$ denoted $d_s(\xi)\ultrato(\xi)_{\nu_s}$.

    The idea now is that we define an (ultra)filter $\mu$ on $S$ such that $(a_s)_{s:\mu}$ gives the ultrafamily from whom $p$ is the ultraretract. We already have maps $(r_s : a_s \ultrato (p)_{\nu_s})$, we need to define some map $i : p\ultrato (a_s)_{s:\mu}$, for that we will use the prebound.

    There is a function $U_p \to \int_{s:\mu}U_{a_s}, \xi \mapsto (d_s(\xi))_{s:\mu}$, let's suppose it comes from a map $i : p\ultrato (a_s)_{s:\mu}$ (we will construct the ultrafilter so that it is the case), then one the composite $U_p \to \int_{s:\mu}U_{a_s} \stackrel{{r_s}_*}\to\int_{s:\mu}(U_p)^{\nu_s}$ can easily be checked to be the diagonal, and so as $U$ is faithful \colorier{(localic maps are faithful on points)} then we get $(r_s)\circ i$ is the thickening of the identity as wanted.

    Only remains to show that we can find an ultrafilter $\mu$ on $S$ such that $U_p \to \int_{s:\mu}U_{a_s}, \xi \mapsto (d_s(\xi))_{s:\mu}$ is in the image of $U$. By the lemma it is sufficient (and necessary) to show that for any $V\subseteq U^n$ and $(\xi_1,\cdots,\xi_n)\in V_p$, then for $\mu$-all $s$, we have $(d_s(\xi_1),\cdots,(d_s(\xi_n))\in V_{a_s}$,
    
    \ie that for all $V\subseteq U^n$ and $(\xi_1,\cdots,\xi_n)\in V_p$ the subsets $S_{V,\overline{\xi}}$ are $\mu$-large
    \[S_{V,\overline{\xi}} \coloneqq \{s : S \ | \ (d_s(\xi_1),\cdots,(d_s(\xi_n))\in V_{a_s}\}.\]

    This family of subsets form a filter basis (as $S_{V,\overline{\xi}}\cap S_{W,\overline{\zeta}} = S_{V\times W, \overline{\xi}.\overline{\zeta}}$), only remains to show that it is a proper filter, \ie that alls theses subsets are non-empty.

    Showing that these subsets are not empty is the technical part of the proof, here we use the hypothesis and we see the need of the thickening $\nu_s$ in the ultra-retract! We fix $V\subseteq U^n$ and $(\xi_1,\cdots,\xi_n)\in V_p$, inspired by the topological case we consider $W\coloneqq V\setminus\cl(\{\overline{\xi}\})$ an open sub-vult of $V\subseteq U^n$ \ie $W$ contains all the $\overline{\zeta}\in V$ such that there are no ultraarow $\overline{\zeta}\ultrato (\overline{\xi})_{\lambda}$ in $V$. As open sub-vult are étale, this form a sheaf over $\X$, but $\overline{\xi}\in V_p\setminus W_p$ so \colorier{using the hypothesis} there exists some $a\in A$ such that $W_a\subsetneq V_a$, \ie there exists $\overline{\zeta}\in V_a$ with maps $\overline{\zeta}\ultrato (\overline{\xi})_{\lambda}$ as wanted. 
\end{proof}

\begin{proof}[preuve alternative]
    We still consider $U\to\X$ étale corresponding to a prebound (and we will denote by $U^n$ the product taken over $\X$). We suppose $p$ is a point of $\E_\res{A}$ and we want to show that $p$ is an ultraretract of points in $A$.
    
Let $I$ be a finite subset of $U_p$, we chose an order $(\xi_1,\cdots,\xi_n) \in (U_p)^n$ \errolnote{Pourquoi un sous-ensemble fini de $U_p$? et les $\xi_i$ sont-ils les éléments de $I$?}\gabrielnote{on demande fini pour se ramener aux $U^n$, et oui c'est juste un choix d'ordre de I, on peut surement travailler avec des listes finies au lieux de finite subsets si tu préfères}. If we consider the topological reflection of $U^n$, we notice that if two opens of $U^n$ coincide on $A$ they coincide on $p$ too \errolnote{Les ouverts de $U^n$ ce sont juste certains faisceaux non? et aussi c'est vrai pour tout faisceau que coincider sur $A$ implique coincider sur $p$ (c'est une façon équivalente de dire que $p$ appartient à $\E_\res{A}$)}\gabrielnote{oui c'est les faisceaux qui sont sous-obj de $U^n$, et oui je suis d'accord sur le deuxième partie}. So the topological case \errolnote{Tu utilises le "topological case" sans dire c'est quoi, c'est pas très clair comment tu l'appliques ici à $U_n$, ni que la conclusion de ce cas satisfait la condition de coherence en bas. Tu as quelque chose de précis en tête clairement mais explicite le}\gabrielnote{j'aplique le résultat à la reflection topologique de $U^n$, l'énoncé du dessus dit exactement que $(\xi_1,\dots,\xi_n)$ appartient à la sobrification (dans la reflection topo de $U^n$) des points de $U^n$ qui vivent au-dessus de $A$, donc que $(\xi_1,\dots,\xi_n)$ est un ultraretract-topologique (ie dans la reflection topo de $U^n$ donc sans coh à priori!) de points de $U^n$ qui vivent au dessus de $A$. Si on écrit c'est quoi un ultraretract dans la reflection topologique=posétal de $U^n$ on obtient exactement ce qui suit } \colorier{(and knowing that the vultraposet of the topological reflection is the posetal reflection!) } gives that $(\xi_1,\cdots,\xi_n)\in (U^n)_p$ is a topological ultraretract of elements in $U^n$ that are above $A$, \ie
    
    there are ultramaps $i^I : p\ultrato(a^I_t)_{t:\nu_I}$ with the $a^I_t$ in $A$, and $(r^I_{t} : a^I_t \ultrato (p)_{\lambda^I_{t}})_{t:\nu_I}$ such that 
    \[((r^I_{t})_{t:\nu_I}\circ i^I)_*(\xi) = (\xi)_{\sum_{t:\nu_I}\lambda^I_{t}} \text{ for $\xi\in I$}.\]

    
    We can merge them by taking $\theta$ a cofinal (ultra)filter on the poset of finite subsets of $U_p$. 
    We get $i : p\ultrato(p)_{I:\theta}\ultrato(a^I_t)_{\sum_{I:\theta}t:\nu_I}$ and we already have the $(r^I_t : a^I_t \ultrato (p)_{\lambda^I_{t}})$. Only remains to show that the composite gives the diagonal, as the prebound is localic \colorier{(and localic are faithful on points!) } it is enough to check that for any $\xi\in U_p$ we have $((r^I_t)_{(I,t)}\circ i)_*(\xi) = (\xi)_{(I:\theta)\otimes(t:\nu_I)\otimes\lambda^I_t}$ \errolnote{il y a une difference entre $\otimes$ et $\Sigma$?} \gabrielnote{non, c'est juste deux notations, à la théorie des types}. But as $\theta$ is cofinal $i_*(\xi) = (i^I)_*(\xi)_{\theta}$ can be computed on the $I$ that contains $\xi$, and so we can use the above formula to conclude.
    \errolnote{Alors j'ai pas où tu utilises qu'on a un ensemble d'ultrafleche conservatives, ni vraiment o` apparaissent les questions de taille. Je suppose que l'hypothèse que $p$ appartient à $\E_\res{A}$ est utilisée dans le "cas topologique" mais re ce que j'ai écris au dessus pour ça.}\gabrielnote{Les notions de tailles n'apparaisent plus dans cette preuve (elles sont cachées dans l'existence de la prebound)}
\end{proof}

\colorier{
\begin{proof}
    We suppose $p$ belongs to $\E_{\res{X}}$ and we fix $U\to\X$ an étale map of v-ult corresponding to a prebound of the topos $\E$.

    In the topological case the indexing set of the ultrafilter can be chosen to be the set of points in the trace of the closure of $p$ on $X$, categorifying this set corresponds to $d\ultrato(a)_{\nu}$ with $a$ above $p$ and $d$ above a point in $X$, but in the categorical case there is not a canonical choice for $a$, so we need to consider such ultraarrows for $a$ varying:
    %\[S=\{s = \{b_1\ultrato(a_1)_{\nu},\cdots,b_n\ultrato(a_n)_{\nu}\} \ | \ \text{ all send to the same atomic $u : x\ultrato (p)_{\nu}$ with $x\in X$}\}\]
    \[\hspace{-30pt}S\coloneqq\{s : (a : U_p)\rightharpoonup\{d_a\stackrel{s_a}\ultrato(a)_{\nu_s}\} \ | \
    \colorier{(\dom(s) \text{ finite ?})}\text{ all $s_a$ are lifts of some common $u_s : x_s\stackrel{at}\ultrato (p)_{\nu_s}$ with $x_s\in X$} \} \]
    $S$ is small as we ask $u_s$ to be atomic, and by the unicity of the lifting $e_a$ is determined by $d_a$.
    
    Let fix some notation: we denote $d_a : S\rightharpoonup U$ the partial map sending $s\in S$ with $a\in\dom(s)$ to $d_a$; for finite sequences $\overline{a} = (a_1,\cdots,a_n)$, we denote $d_{\overline{a}} : S\rightharpoonup U^n$ for the product of the $d_{a_i}$; and for $V\subseteq U^n$, we denote $S_{V,\overline{a}} \coloneqq \{s\in S \ |\ d_{\overline{a}}(s) \in V\}$.

    In the topological case, the main ingredient of the proof is to show that every open containing $p$ intersects the trace of the closure of $p$ on $X$, this categorifies into 
    \begin{nclaim}
        The sets $S_{V,\overline{a}}$ are non-empty.
    \end{nclaim}
    \begin{proof}[proof of the claim]
        Take $V\subseteq U^n$ and $\overline{a}\in V_p$. Note that $V\subseteq U^n$ is still an étale space over $\X$ and consider the full sub-v-ult of $V$
        \[W\coloneqq V{\setminus\{\overline{b} \ | \ \exists \overline{b}\ultrato(\overline{a})_{\mu}\}}\]
        and note that it is an étale space over $\X$ (the unique lifting property is still satsified). As $\overline{a}\in V_p\setminus W_p$ we can find some $x\in X$ such that $W_x\subsetneq V_x$ \ie there is $(b_1,\cdots,b_n)\in V_x{\setminus W_x}$. This means exactly that $b_i\in V_x$ and that there is some $(b_i\ultrato (a_i)_{\mu})$ that live above the same $x\ultrato(p)_{\mu}$. We can then maximally defatten this ultraarow to make it atomic.
    \end{proof}

    We notice that $S_{V,\overline{a}}\cap S_{W,\overline{b}} = S_{V\times W, \overline{a}.\overline{b}}$, and so the $(S_{V,\overline{a}})$ can be extended to an ultrafilter $\mu\in\beta(S)$. We get $(x_s)_{s:\mu}$ an ultrafamily of points of $X$ and we can define $v : p\ultrato(x_s)_{s:\mu}$ using the prebound:
    
    $\alpha : \overline{a} \in U_p\mapsto (d_a(s))_{s:\mu} \in \int_{s:\mu}U_{x_s}$
    that is well-defined as $\dom(d_a) = S_{U,a}$
    
    for all $V\subseteq U^n$, $\alpha^n$ restricts to $V_p \to \int_{s:\mu}V_{x_s}$ as for $s\in S_{V,\overline{a}}$, $d_{\overline{a}}(s)\in V_{x_s}$
    
    hence we get a natural transformation $\Nat_{E:\E}(E_p,\int_{s:\mu}E_{x_s})$, defining thus a $v\in\X(p,(x_s)_{s:\mu})$.

    Finally, for $s:\mu$ we have $u_s : x_s\ultrato (p)_{\nu_s}$, and we want to check that $(u_s)\circ v$ is the ultradiagonal. As $U:\X\to\Set$ is faithful, it is enough to check it on any $a\in U$ and this comes from the definition of $S$.
\end{proof}}

\begin{defn}
    Let $\X$ be a v-ult. An ultraidempotent in $\X$ is given by
    \begin{enumerate}
        \item objects $(a_s)_{s:\mu}$
        \item ultrafilters $(\nu_s)_{s:\mu}$
        \item ultraarrows $(\gamma_{s_0} : a_{s_0}\ultrato(a_s)_{\nu_{s_0}\otimes(s:\mu)})_{s_0:\mu}$.
        \item such that for $s_0:\mu$
        \[a_{s_0}\stackrel{\gamma_{s_0}}\ultrato(a_s)_{\nu_{s_0}\otimes(s:\mu)} \stackrel{(\gamma_s)}\ultrato((a_t)_{\nu_s\otimes(t:\mu)})_{\nu_{s_0}\otimes(s:\mu)}\]
        corresponds to 
        \[a_{s_0}\stackrel{\gamma_s}\ultrato(a_t)_{\nu_{s_0}\otimes(t:\mu)}\]
        trough the fattening given by \[\nu_{s_0}\otimes(s:\mu)\otimes\nu_s\otimes(t:\mu) \to \nu_{s_0}\otimes(t:\mu)\]
    \end{enumerate}
\end{defn}

\begin{defn}
    The splitting of such an ultraidempotent is given by
    \begin{enumerate}
        \item $(u_s : a_s\ultrato (p)_{\nu_s})$
        \item $v : p\ultrato(a_s)_{s:\mu}$
        \item such that
        \[p\stackrel{v}\ultrato (a_s)_{s:\mu}\stackrel{(u_s)}\ultrato ((p)_{\nu_s})_{s:\mu}\]
        is the diagonal (the fattening of the identity),
        \item and
        \[a_{s_0}\stackrel{u_{s_0}}\ultrato (p)_{\nu_{s_0}}\stackrel{(v)}\ultrato ((a_s)_{s:\mu})_{\nu_{s_0}}\]
        is equal to $\gamma_{s_0}$ for $s_0:\mu$.
    \end{enumerate}
\end{defn}

\begin{ex}
    For a single object $a$ (\ie $\mu = \ptuf$) and all $\nu_s=\ptuf$, we recover the usual notion of idempotent and splitting of idempotent
\end{ex}

\begin{ex}
    For $(u_s : a_s\ultrato (p)_{\nu_s})$ and $v : p\ultrato(a_s)_{s:\mu}$ such that
    $p\stackrel{v}\ultrato (a_s)_{s:\mu}\stackrel{(u_s)}\ultrato ((p)_{\nu_s})_{s:\mu}$
    is the diagonal, the composite 
    \[\gamma_s : a_{s_0}\stackrel{u_{s_0}}\ultrato (p)_{\nu_{s_0}}\stackrel{(v)}\ultrato ((a_s)_{s:\mu})_{\nu_{s_0}}\]
    give an ultraidempotent on $(a_s)_{s:\mu}$ and $p$ is a spliting of it.
\end{ex}

\begin{ex}
    In an ultracategory, $A$ is a splitting of $A^{\nu}$ with $u : A^{\nu}\ultrato (A)_{\nu}$ the identity of $A^{\nu}$ and $v : A\to A^{\nu}$ the diagonal. 
\end{ex}

\begin{ex}
    The v-ult $\Set$ has splitting of its ultra-idempotent. Explicitly, for a $(A_s)_{s:\mu}$ and $\gamma_{s_0} : A_{s_0}\to(\int_{s:\mu}(A_s))^{\nu_{s_0}}$ an ultra-idempotent in $\Set$, the splitting is given by
    \[P\coloneqq\{(a_s)_{s:\mu} \ |\ (\forall s:\mu) \gamma_s(a_s) = ((a_s)_{s:\mu})_{t_{\nu_{s_0}}} \}\]
    with $u_s : A_s \to P^{\nu_{s_0}} \ ;\ a\mapsto \gamma_s(a)$ and $v : P\to\int_{s:\mu}A_s$ the inclusion.    
\end{ex}

\begin{ex}
    Let $(a_i)_{i:I}$ be a filtered diagram in $\pt(\X)$ and $\mu\in\beta(I)$ a cofinal ultrafilter (\ie extending the filter of upsets). For each $i:I$ we define $\gamma_i : a_i \ultrato (a_i)_{s:\mu}\to(a_s)_{s:\mu}$, then $(\gamma_s)$ is an ultra-idempotent.

    Moreover, if $p$ is a colimit of $(a_i)_{i:I}$ satisfying the universal property for ultraarows, then then $(\gamma_i)$ induce a $v : p\ultrato(a_s)_{s:\mu}$ and taking the $(u_s:a_s\to p)$ being the map of the universal cocones, we get that $p$ is a spliting of this ultra-idempotent.

    Conversely, given $p$ a spliting, then we get a cocone $(u_s : a_s\to p)$ and given ultraarrows $(a_s\ultrato (y_t)_{t:\nu})$ we can define $(p \ultrato (a_s)_{s:\mu}\ultrato ((y_t)_{t:\nu})_{s:\mu})$, and if the ultrad condition is satisfied, we can reduce it to an ultraarow $p\ultrato(y_t)_{t:\nu}$.

    Hence (modulo the ultrad condition) giving a splitting of $(a_s,\gamma_s)$ is equivalent to a giving an ultracolimit of $(a_i)_{i:I}$ (\ie a colimit with the universal property for ultraarows and not only arrows).
\end{ex}

\begin{prop}
    Let $(p,v,u_s)$ be a splitting of $(a_s,\nu_s,\gamma_s)$ as above. Let $f : p\to(y_l)_{l:\lambda}$ an ultraarow, precomposing with the $u_s$ gives ultraarows $(f_s : a_s\ultrato((y_l)_{l:\lambda})_{\nu_s})$ that are compatible in the sense that for $s_0:\mu$
    \[a_{s_0} \stackrel{\gamma_{s_0}}\ultrato ((a_s)_{s:\mu})_{\nu_{s_0}} \stackrel{(f_s)}\ultrato ((((y_l)_{l:\lambda})_{\nu_s})_{s:\mu})_{\nu_{s_0}}\]
    corresponds to 
    \[a_{s_0}\stackrel{f_{s_0}}\ultrato((y_l)_{l:\lambda})_{\nu_{s_0}})\]
    trough the fattening given by $\nu_{s_0}\otimes(s:\mu)\otimes\nu_s\otimes\lambda \to \nu_{s_0}\otimes\lambda$.

    Conversely, given ultraarows $(f_s : a_s\ultrato((y_l)_{l:\lambda})_{\nu_s})$ that are compatible in the above way, we can precompose by $v$
    \[p\stackrel{v}\ultrato(a_s)_{s:\mu}\stackrel{(f_s)}\ultrato(((y_l)_{l:\lambda})_{\nu_s})_{s:\mu}\]
    and de-fatten into some $f : p\ultrato(y_l)_{l:\lambda}$
\end{prop}

\begin{proof}
    We use the ultrad condition, we show that fattening this composite by
    \[(s:\mu)\otimes\nu_s\otimes(t:\mu)\otimes\nu_t\otimes\lambda\to(s:\mu)\otimes\nu_s\otimes\lambda\] is the same than precomposing it with the diagonal $p\ultrato((p)_{\nu_s})_{s:\mu}$. But fattening the composite can be reduced to fattening the $(f_s)$ part only, and by the compatibility condition the $(f_s)$ it is the same as 
    \[p\stackrel{v}\ultrato(a_s)_{s:\mu}\stackrel{(\gamma_s)}\ultrato(((a_t)_{t:\mu})_{\nu_s})_{s:\mu}\stackrel{f_t}\ultrato(((((y_l)_{l:\lambda})_{\nu_t})_{t:\mu})_{\nu_s})_{s:\mu}\] 
    and using the splitting equation $v$ followed by the $(\gamma_s)$ is the ultradiagonal followed by $(v)$.
\end{proof}

\begin{prop}
    Moreover these two operations establish a bijective correspondence between compatible families $(f_s : a_s\ultrato((y_l)_{l_\lambda})_{\nu_s})$ and ultraarows $f:p\ultrato(y_l)_{l:\lambda}$.
\end{prop}
\begin{proof}
    Follows from the splitting equations. In both directions, we use the unicity of the de-fattening when it exists.
\end{proof}

\begin{cor}
    For $(y_l)_{l:\lambda}$ in $\bb(\X)$, the homset $\Hom(p,(y_l))$ is in bijection with compatible families $\Comp((a_s,\gamma_s), (y_l))$; moreover, this bijection is natural in $(y_l)\in\bb(\X)$.
\end{cor}

This is the colimit universal property, there is also a limit one: maps $x\to p$ correspond to $x\ultrato(a_s)_{s:\mu}$ that are stable by postcomposition with the $\gamma_s$ (and similarly for $x\to(p)_{\kappa}$). This universal property is easier to prove, and gives also the uniqueness of the splitting. 

Conjecture: Let $X$ be a subclass of points the following are equivalent
\begin{enumerate}
    \item 
\end{enumerate}


\section{Faithful functors preserve local soberness}

On commence par montrer la proposition suivante (dans le cas $U = 1$ on obtient que la topologie à la Sam coincide avec la reflection posetale)
\begin{prop}
    Soit $\X$ une petite ultrad, $U : \X\to\Set$. Alors une set-fonction $f\in\Set(U(p),U(q))$ est dans l'image de $U$ (up to thick) ssi 
    \[\forall A \subseteq U^n, A(p)\ni(a_1,\dots, a_n)\to A(q)\ni(f(a_1),\dots, f(a_n))\]
\end{prop}

\begin{proof}
    Le sens direct est clair, on suppose donc que $f$ vérifie la condition et on veut montrer que $f$ s'épaissie en une map de la forme $U(e)$ pour un certain $e : p \ultrato(q)_{\mu}$.

    \begin{nclaim}
        Pour toute partie finie $I\subseteq U_p$, on peut trouver $e_I : p \ultrato(q)_{\mu_I}$ telle que $U(e_I)(a) = f(a)$ pour tout $a\in I$.
    \end{nclaim}

    \begin{proof}[proof of the claim]
        Pour fixer les idées on suppose $I=\{a,b\}$ avec $a\neq b$. Alors tout ouvert de $\U\times_{\X}\U$ (ou $\U$ désigne l'espace étale associée à $U$), contenant $(a,b)$ contient aussi $(fa,fb)$. Donc (modulo que $\U\times_{\X}\U$ est une ultrad et donc que sa reflection posetal coincide avec sa 0-topologie), on obtient une ultraflèche $(a,b)\ultrato_{\mu_I}(fa,fb)$ dans $\U\times_{\X}\U$, ce qui correspond à deux ultraflèches $a\ultrato(fa)_{\mu_I}$ et $b\ultrato(fb)_{\mu_I}$ qui vivent au-dessus d'une même ultraflèche $e_I : p\ultrato(q)_{\mu_I}$.
    \end{proof}
    
    On peut alors rassembler les $e_I$ en une ultraflèche $e : p \ultrato (p)_{\nu} \stackrel{(e_I)}\ultrato(q)_{\sum_{\nu}\mu_I}$, ou on a choisi $\nu$ un ultrafiltre sur les parties finies de $I\subseteq \U_p$ qui étend l'ordre (\ie l'ensemble des parties finies contenant un certain $a\in U_p$ est large).

    On prend un $a\in U(p)$ et calcul $U(e)(a) = (e_I(a))_{I:\nu}$, comme pour presque tout $I:\nu$ on a $e_I(a) = (f(a))_{\mu_I}$ on obtient $U(e)(a) = (f(a))_{\sum_{\nu}\mu_I}$, \ie $U(e)(a) = (f(a))_{\mu}$ pour tout $a\in U(p)$ avec $\mu := \sum_{\nu}\mu_I$
    
    Donc $U(e)$ correspond à l'épaississement $U_p\stackrel{f}\to U_q\hookrightarrow(U_q)^{\mu}$ comme voulu.
\end{proof}


\begin{prop}
    Si de plus $U : \X\to\Set$ est fidèle, alors $X$ est localement sobre, \ie que $\X\to\pt(\Sh(X))$ est ff.
\end{prop}
\begin{proof}
    On note $\E := \Sh(\X)$.
    
    On voir facilement que $sob:\X\to\pt(\E)$ est faithful: si $\alpha,\beta:p\ultrato(...)$ sont parallèles tels que $sob(\alpha) = sob(\alpha) : \Nat_{E:\E}(E_p,...)$, en évaluant ces deux transfo nat en $U:\E$ on obtient $U(\alpha) = U(\beta)$ et on conclut en utilisant $U$ fidèle.

    Plus simplement, décomposer $U$ en $U : \X\to\pt(\E)\stackrel{\pt\sh U}\to\Set$ qui est faithful.
    
    Il nous reste donc à montrer qu'il est full up to thickening.

    On prend une transformation naturelle $\alpha : \Nat_{E:\E}(E_p,E_q) = \pt(\E)(p,q)$. On obtient une map $\alpha_U : U_p\to U_q$ qui vérifie la proposition précédente, et donc un $e : p\ultrato(q)_{\mu}$ tel que $\alpha'_U = U(e)$ (où $\alpha'$ dénote l'épaississement de $\alpha$ par $\mu$).
    \ie on a $\alpha'_U = sob(e)_U$.
    
    Pour en déduire que $\alpha' = sob(e)$, il suffirait de montrer que $U$ est une préborne sur $\E$, et donc une tel transfo nat est tjs déterminée par sa valeure en $U$...
\end{proof}


\section{The wisdom of the prebound}

We know that $\X\to\pt(\E)$ full up to thick implies $\sh(\X)\to\E$ hyperconnected.

\begin{defn}
    If $\X$ is a v-ult, its \emph{topological reflection} its topological reflection is $\X\to P$ where $P$ has the same objects than $X$ and there is a unique ultraarow $a\ultrato(b_s)$ in $P$ iff there exists one $a\ultrato(b_{\alpha(t)})$ in $\X$ for some thickening. 
\end{defn}
\begin{rmk}
    $P$ is a (potentially large) topological space, and is a topological space if $\X$ is accessible. Moreover it is the universal topological space with such a map. 
\end{rmk}
\begin{lem}
    If $\X$ v-ult, and $\X\to P$ its topological reflection. Then 
    \[\X\to\pt(\sh(\X))\to\pt(Loc(\sh(\X))) = \pt(open(\X))\]
    gives a map $P \to \pt(open(\X))$ that is the soberification of $P$. 
\end{lem}
\begin{proof}
    Go to the sheaf: the composite
    \[\sh(\X)\to\sh(P)\to Loc(\sh(\X))\]
    is hyperconnected and $\sh(P)\to Loc(\sh(\X))$ is localic.

    Hence it is enough to show $\sh(\X)\to\sh(P)$ hyperconnected to deduce by unicity of the hyperconnected/localic factorization. This is enough by the direction full implies hyperconnected.  
\end{proof}

\begin{defn}
    For $A : \X\to\Set$ an ultrasheaf, we denote $\ty(A) \in \TopSp$ the topological reflection of the étale space $\A\to\X$ associated to $\A$. 
\end{defn}
\begin{rmk}
    the soberification of $\ty(A)$ can be described also as the points of the local $Sub_{\sh(\X)}(A)$ of subobjects of $A\in\sh(\X)$ or equivalently as the topological space of points of the localic reflection of $\sh(\X)/A$.
\end{rmk}

\begin{Thm}
    Let $\X$ a v-ultracat, $U : \X\to\Set$ an ultrasheaf on $\X$, $p,(q_s)_{s:\mu}$ points of $\X$, $\U\stackrel{\text{ét}}\to\X$ the étale space associated to $U$ and $\ty(U^n)$ the topological reflection of  $\U\times_\X\times_\X \cdots \U$.

    We moreover suppose $U(p)$ non-empty.
    
    For $f : U(p)\to \int U(q_s)$ the following are equivalents:
    \begin{enumerate}
        \item $f$ is in the image of $U$ \ie there is $\alpha : \nu\to\mu$ and $e : p \ultrato (q_{\alpha(t)})_{t:\nu}$ such that $U(e) = \alpha^*f$.
        \item for $\overline{a}\in U(p)^n$ (so that $f(\overline{a})\in \int (U(q_s)^n) = (\int (U(q_s))^n$) 
        \[\forall A\subseteq U^n,\ \overline{a}\in A(p) \Rightarrow f(\overline{a})\in \int A(q_s)\] 
        \ie $f^n : U(p)^n\to\int_{s:\mu}U(q_s)^n$ restricts to $A(p)\to\int_{s:\mu}A(q_s)$.
        
        \item or equivalently
        \[\overline{a}\vuleq{} (f(\overline{a})) \text{ in } \ty(U^n).\]

    \end{enumerate}
    If moreover $U$ is faithful, we can add
    \begin{enumerate}
        \item[(iv)] $f$ is in the strict image of $U$ \ie $f = U(e)$ for some $e : p\ultrato(q_s)$.
    \end{enumerate}
\end{Thm}
Note that the lemma above corresponds to the case $U = 1$. Logically the lemma says that if any geometric formula satsfied by $M$ is also satisfied by $N$ then there is a morphism from $M$ to an ultrapower of $N$.

Logically $U$ is tough of as a sort of the models of $\X$ (not necessarily a generating one), and this proposition says that a set-function defined on the underlying set specified by $U$, respects all the formulas with parameters in $U$ iff it can be thicken so that it arises from a morphism of models. It is actually be also seen as a special case of the above by adding constant to the theory.

Maybe separated the prop with first the case $U(p)$ finite?

We prove (iii) implies (i) the other implications are not too hard.
\begin{proof}[Proof of (iii) implies (i)]
     We suppose $f$ satisfies (iii), and want to construct an $e : p \ultrato (q_{\alpha(t)})$ for some $\alpha : \nu\to\mu$ such that $U(e) = \alpha^*f$.

     First, we construct for any $I\subseteq U(p)$ finite subset, an ultraarow $e_I : p \ultrato(q_{\alpha_I(t)})_{t:\nu_I}$ such that $U(e_I)(a) = \alpha_I^*f(a)$ for $a\in I$. Let say $I := \{a_1,\dots,a_n\}$, the hypothesis plus the lemma above gives an ultraarow $(a_1,\cdots,a_n)\ultrato_{\nu_I}(f(a_1),\cdots,f(a_n))$ with $\alpha_I : \nu_I\to\mu$ in $\U\times_\X\cdots\times_\X\U$, and so $n$ ultraarows $(a_i\ultrato_{\nu_I}(f(a_i)))$ that live above the same $e_I : p\ultrato(q_{\alpha_I(t)})_{t:\nu_I}$.

     Then by choosing a cofinal ultrafilter $\lambda$ on the poset of finite subsets of $U(p)$, so that we can merge the $e_I$ into an ultraarow $e : p \ultrato (q_{\alpha_I(t)})_{\sum_{I:\lambda}(t:\nu_I)}$. And for any $a \in U(p)$, one can compute $U(e)(a) = \alpha^*f(a)$ where $\alpha : \sum_{I:\lambda}(t:\nu_I) \to \mu$ is given by the sum of the $\alpha_I$.
\end{proof}

\begin{rmk}
    This theorem doesn't use the reconstruction theorem (to verify!).
\end{rmk}

\begin{rmk}
    The lemma \ref{Lemprebound} gives another proof (more elementary?) of the special case of $\X$ sober and $U$ prebound.
\end{rmk}

The crucial corollary is the following.
\begin{cor}
    Let $U : \X\to\Set$ an ultrasheaf, and consider the hyperconnected/localic factorization $\Sh(\X)\twoheadrightarrow\E\to\SetO$. Then $g : \X\to\pt(\E)$ is the full up to thickening.
\end{cor}
\begin{proof}
    We consider $p,(q_s)$ in $\X$ and we take $h : g(p)\ultrato(g(q_s))_{s:\mu}$. First note that $V : \pt(\E)\to\Set$ is faithful as points of a localic geometric morphism, $h$ induces a map $V(h) : U(p)\to\int_{s:\mu}U(q_s)$.
    
    It is enough to show that satisfies the above hypothesis, so that we get $k : p \ultrato(q_{\alpha(t)})_{t:\nu}$ with $U(k) = \alpha^*V(h)$, \ie $V(g(k)) = V(\alpha^*h)$ and we conclude that $g(k) = \alpha^*h$ by faithfulness of $V$.

    Let's show that $V(h)$ satisfies the hypothesis of the theorem above. There is a geometric morphism $\Sh(\X)/(U^n) \to \E/(V^n)$ that is hyperconnected as a pullback of an hyperconnected morphism. So that the localic relfection given by the morphism of locals $Sub_{\Sh(\X)}(U^n) \to Sub_{\E}(V^n)$ is an equivalence and taking the points we get an equivalence of topological spaces $\ty(U^n)\to\ty(V^n)$.
\end{proof}

\begin{cor}
    \begin{enumerate}
        \item For $f : \sh(\X)\twoheadrightarrow\E$ hyperconnected, $\X\to\pt(\E)$ is full up to thickening.
        \item If there exists $U : \X\to\Set$ faithful with $\X$ ultrad, then $\X$ is locally sober and $\sh(X)\to\Set$ is localic.
        \item More generally if $U : \X\to\pt(\E)$ is faithful with $\X$ ultrad, then $\X$ is locally sober and $\sh(X)\to\E$ is localic.
    \end{enumerate}
\end{cor}
\begin{proof}
    \begin{enumerate}
        \item Chose a prebound of $\E$ and apply the above.
        \item Using that $\X$ ultrad implies that faithful plus full up to fhick gives $\X$ fullyfaithful into $\pt(\E)$ sober, we get $\X$ locally sober, and moreover $\sh(\X)\to\E$ embedding.
    \end{enumerate}
    
\end{proof}
